\chapter{Аналитическая часть}

\section{Графы и связность}

Граф $G$ определяется множеством вершин~(\ref{eq:vert_set}) и множеством рёбер~(\ref{eq:edge_set}):
\begin{equation}
	V = \{ v_1, v_2, \ldots, v_n \},
	\label{eq:vert_set}
\end{equation}
\begin{equation}
	E = \{ (v_i, v_j) \mid v_i, v_j \in V \},
	\label{eq:edge_set}
\end{equation}

Таким образом, граф задаётся парой множеств~(\ref{eq:graph_pair_def}):
\begin{equation}
	G = (V, E).
	\label{eq:graph_pair_def}
\end{equation}

Путь в графе представляет собой последовательность вершин, каждые две соседние из которых соединены ребром.  
Граф называется связным, если для любых двух различных вершин существует путь, соединяющий их.  
Если хотя бы до одной вершины нельзя добраться из другой, граф считается несвязным~\cite{aho}.

\section{Удаление вершины}

В соответствии с вариантом задания требуется определить, сохраняет ли граф связность после удаления вершины $k$.  
Удаление вершины означает исключение её из множества $V$ и удаление всех рёбер, в которых она участвовала.  
Полученный граф обозначается следующим образом~(\ref{eq:graph_after_del}):

\begin{equation}
	G' = (V \setminus \{k\}, E').
	\label{eq:graph_after_del}
\end{equation}

Граф $G'$ считается связным, если оставшиеся вершины продолжают образовывать единую компоненту связности.

\section{Последовательный алгоритм проверки связности}

Для определения связности используется обход графа в ширину (Breadth-First Search, BFS).  
Этот метод позволяет установить, какие вершины достижимы из выбранной стартовой вершины.

Последовательный алгоритм включает следующие этапы:
\begin{enumerate}
	\item считывается исходный граф $G$.
	\item выполняется удаление вершины $k$ и всех рёбер, связанных с ней, формируя граф $G'$.
	\item выбирается стартовая вершина $s \in V \setminus \{k\}$, имеющая хотя бы одно ребро.
	\item из вершины $s$ выполняется обход BFS:
	\begin{itemize}
		\item создаётся очередь вершин;
		\item каждая посещённая вершина отмечается в массиве признаков;
		\item в очередь помещаются все её соседи.
	\end{itemize}
	\item после завершения обхода проверяется, посещены ли все вершины графа $G'$.
\end{enumerate}

Если посещены все вершины, граф считается связным.

\section{Параллельный алгоритм проверки связности}

Для ускорения вычислений используется параллельная модификация обхода BFS.  
Основная идея заключается в том, что несколько потоков могут одновременно извлекать вершины из общей очереди и обрабатывать их соседей.

Параллельный алгоритм включает следующие шаги:
\begin{enumerate}
	\item определяется стартовая вершина $s$ после удаления вершины $k$.
	\item создаётся общая очередь, содержащая стартовую вершину.
	\item запускается заданное число рабочих потоков.
	\item каждый поток:
	\begin{itemize}
		\item извлекает вершину из очереди;
		\item обрабатывает всех её соседей;
		\item помещает непосещённых соседей в очередь.
	\end{itemize}
	\item работа потоков завершается, когда очередь пуста и все потоки находятся в состоянии ожидания.
\end{enumerate}

Такой подход позволяет обрабатывать различные части графа параллельно, уменьшая время выполнения на больших входных данных.

\section{Основные понятия параллельного программирования}

При реализации параллельной версии алгоритма используются базовые понятия.

Процесс представляет собой исполняемую программу, обладающую собственным адресным пространством и ресурсами.  
Процессы изолированы друг от друга, что обеспечивает их независимость~\cite{zaharov}.

Поток является единицей выполнения внутри процесса.  
Потоки разделяют память процесса и могут выполняться параллельно, обрабатывая одни и те же данные~\cite{stolyarov,zaharov}.

Мьютекс служит для обеспечения взаимного исключения.  
Он позволяет организовать доступ к участку программы таким образом, чтобы в определённый момент времени этот участок мог выполняться только одним потоком.  
Использование мьютекса предотвращает гонки данных при работе с общей памятью~\cite{stolyarov}.

Синхронизация потоков представляет собой набор методов согласования действий потоков, необходимых при совместном доступе к данным.  
Она обеспечивает согласованность состояния программы и предотвращает нарушения целостности данных~\cite{zaharov}.

\section*{Вывод}

В аналитической части были рассмотрены основные определения графов и связности, а также сформулированы последовательный и параллельный алгоритмы проверки связности графа после удаления вершины.  
Последовательная версия основана на обходе графа в ширину.  
Параллельная версия распределяет обработку между несколькими потоками и требует использования механизмов синхронизации для обеспечения корректного доступа к общим данным.
